\makeabstract
    {
    Dependence of the \ensuremath{\pi}{\textasciicircum}+-e{\textasciicircum}+ Miss Distance on the Active TARget{\textquoteright}s Position Resolution for the PIONEER Experiment.
    }
    {
    Audren Hedges Duroy\textsuperscript{1}, Sean Foster\textsuperscript{1}
    }
    {
    \textsuperscript{1} Department of Physics \& Astronomy, Amherst College, Amherst, MA
    }
    {
    This study investigates the impact of point resolution on the ability for a GenFit fitter to trace positron tracklets back to the pion stopping point from which they decayed for the PIONEER Experiment. Because both Multiple Coulomb Scattering and the PIONEER Active TARget's (ATAR) detector point resolution impact a fit in the reconstruction of pion decay events, this study aims to examine the importance of increased hit point resolution.
    }
    {
    This study examined \ensuremath{\pi}{\textrightarrow}e\ensuremath{\nu} decay events. It employed the PIONEER simulation software, built with Geant4, to collect the true location of ATAR hits, and apply artificial resolutions. The uncertainty in hit position was set to these chosen resolutions, and the true hits were smeared using random numbers within these uncertainties, to simulate measurement resolution. Following a series of cuts to ensure clear event collection, a Kalmann filter was applied, using the GenFit architecture, to extrapolate the track of positron hits back to the true stopping location of the pion in each event. The difference in the extrapolated positron position and the pion stopping point was aggregated across 8,000 events, per chosen resolution, and an analysis of the resulting histogram standard deviations for each resolution was performed.
    }
    {
    As resolution decreased beyond the current true resolution (57.7 microns), a continued downward trend in standard deviations of stopping distances was observed. Indeed, reducing the detector point resolution by a factor of two reduced the miss distance uncertainty by 25\%.
    }
    {
    This study presents compelling evidence that, for \ensuremath{\pi}{\textrightarrow}e\ensuremath{\nu} events in the PIONEER Experiment, there is merit in increased detector resolution. The geometry of the ATAR, through the effects of Multiple Coulomb Scattering, does not appear to outweigh the impact of increased point resolution.
    }

    \newpage
    